\documentclass[11pt,a4paper]{article}

% ============================================================
%  QIMMA -- COURS : Arithmétique dans Z
%  2ème Bac Sciences Mathématiques (A et B)
%  Standard exhaustif : définitions formelles + démonstrations
%  Basé sur templates/qimma-template.tex (charte Qimma)
% ============================================================

\usepackage[french]{babel}
\usepackage[utf8]{inputenc}
\usepackage[T1]{fontenc}
\usepackage{lmodern}
\usepackage[a4paper,margin=2cm,top=2.3cm,bottom=2.6cm]{geometry}
\usepackage{amsmath,amssymb}
\usepackage{xcolor}
\usepackage{graphicx}
\usepackage{tikz}
\usetikzlibrary{shadows,calc}
\usepackage[most]{tcolorbox}
\usepackage{titlesec}
\usepackage{fancyhdr}
\usepackage{enumitem}
\usepackage{pifont}
\usepackage{array}
\usepackage{colortbl}
\usepackage{hyperref}

% ------------------- CHARTE QIMMA -------------------
\definecolor{qteal}{HTML}{0d9488}
\definecolor{qtealdark}{HTML}{134e4a}
\definecolor{qamber}{HTML}{fbbf24}
\definecolor{qambertext}{HTML}{92400e}
\definecolor{ink}{HTML}{1f2937}
\definecolor{inkgray}{HTML}{6b7280}
\definecolor{tealpale}{HTML}{e6f5f3}
\colorlet{colDef}{qteal}
\definecolor{colProp}{HTML}{0f766e}
\definecolor{colEx}{HTML}{b45309}
\definecolor{colRem}{HTML}{9d174d}
\color{ink}

\graphicspath{{../../../../assets/}}
\newcommand{\logofile}{../../../../assets/qimma-logo.pdf}

% ------------------- METADONNEES -------------------
\newcommand{\matiere}{Mathématiques}
\newcommand{\niveau}{2\up{ème} Bac -- Sciences Mathématiques}
\newcommand{\chapitretitre}{Arithmétique dans $\mathbb{Z}$}

% ------------------- EN-TETE / PIED DE PAGE -------------------
\pagestyle{fancy}
\fancyhf{}
\renewcommand{\headrulewidth}{0pt}
\renewcommand{\footrulewidth}{0.5pt}
\fancyfoot[L]{\raisebox{-0.15cm}{\includegraphics[height=0.35cm]{\logofile}}~\small\sffamily\color{inkgray}Qimma}
\fancyfoot[C]{\small\sffamily\color{inkgray}\matiere\ -- \niveau}
\fancyfoot[R]{\small\sffamily\color{qtealdark}\thepage}

% ------------------- TITRES DE SECTION -------------------
\titleformat{\section}
  {\normalfont\sffamily\Large\bfseries\color{qtealdark}}
  {\tikz[baseline=-3.2pt]{\node[circle,fill=qteal,text=white,inner sep=0pt,minimum size=8mm]{\sffamily\bfseries\thesection};}}
  {10pt}{}
  [\vspace{-2mm}{\color{qamber}\titlerule[1.6pt]}\vspace{2mm}]
\titlespacing*{\section}{0pt}{16pt}{10pt}

\titleformat{\subsection}
  {\normalfont\sffamily\large\bfseries\color{qtealdark}}
  {\color{qteal}\ding{212}}{6pt}{}
\titlespacing*{\subsection}{0pt}{12pt}{6pt}

% ------------------- BOITES PEDAGOGIQUES -------------------
\newtcolorbox{fiche}[3][]{
  enhanced, breakable,
  colback=white, colframe=#2,
  boxrule=1pt, arc=2mm,
  top=7mm, left=4mm, right=4mm, bottom=3mm,
  drop shadow={black!25, shadow xshift=1mm, shadow yshift=-1mm},
  attach boxed title to top left={xshift=4mm, yshift=-3mm, yshifttext=-1mm},
  boxed title style={colback=#2, arc=1.5mm, boxrule=0pt, left=2mm, right=2mm},
  coltitle=white, fonttitle=\sffamily\bfseries,
  title={#3}, #1
}
\newenvironment{definition}[1][Définition]
  {\begin{fiche}{colDef}{\ding{110}~~#1}}{\end{fiche}}
\newenvironment{propriete}[1][Propriété]
  {\begin{fiche}{colProp}{\ding{72}~~#1}}{\end{fiche}}
\newenvironment{exemple}[1][Exemple]
  {\begin{fiche}{colEx}{\ding{217}~~#1}}{\end{fiche}}
\newenvironment{remarque}[1][Remarque]
  {\begin{fiche}{colRem}{\ding{43}~~#1}}{\end{fiche}}

% Démonstration (hors boîte, style discret)
\newenvironment{demo}
  {\par\smallskip\noindent{\sffamily\bfseries\color{qtealdark}Démonstration.}\ \itshape}
  {\ \normalfont\hfill$\blacksquare$\par\medskip}

\hypersetup{colorlinks=true, linkcolor=qtealdark, urlcolor=qtealdark}

% ============================================================
\begin{document}

% ------------------- BANNIERE DE TITRE -------------------
\begin{tcolorbox}[
  enhanced, boxrule=0pt, arc=4mm,
  interior style={left color=qtealdark, right color=qteal},
  top=8mm, bottom=8mm, left=10mm, right=32mm,
  drop shadow={black!35, shadow xshift=1.5mm, shadow yshift=-1.5mm},
  before skip=0pt, after skip=9mm,
  overlay={
    \node[anchor=north west] at ([xshift=6mm,yshift=-6mm]frame.north west)
      {\includegraphics[height=1.25cm]{\logofile}};
    \node[anchor=north east, fill=qamber, text=qambertext, rounded corners=2pt,
          inner sep=4pt, font=\sffamily\bfseries\small] at ([xshift=-6mm,yshift=-6mm]frame.north east) {COURS};
  }
]
\hspace{1.6cm}\centering
{\sffamily\bfseries\color{white}\fontsize{23}{27}\selectfont \chapitretitre}\\[3mm]
{\sffamily\color{white!90}\large \matiere\ \textbullet\ \niveau}
\end{tcolorbox}

% ------------------- OBJECTIFS -------------------
\begin{tcolorbox}[
  enhanced, colback=tealpale, colframe=qteal, boxrule=0.8pt, arc=2mm,
  left=5mm, right=5mm, top=3mm, bottom=3mm,
  title={\sffamily\bfseries\color{qtealdark}\ding{51}~~Objectifs du chapitre},
  coltitle=qtealdark, colbacktitle=tealpale, boxrule=0.8pt,
  attach title to upper={\par\vspace{1mm}}
]
\begin{itemize}[leftmargin=6mm, label=\ding{212}, itemsep=1mm]
  \item Maîtriser la divisibilité dans $\mathbb{Z}$ et démontrer ses propriétés fondamentales.
  \item Effectuer et justifier la division euclidienne, calculer PGCD et PPCM, et démontrer leurs propriétés.
  \item Appliquer l'algorithme d'Euclide et le théorème de Bézout, résoudre des équations diophantiennes $ax+by=c$.
  \item Démontrer et utiliser le théorème de Gauss, caractériser les nombres premiers entre eux.
  \item Manipuler les congruences modulo $n$, résoudre des équations et démontrer des critères de divisibilité.
  \item Écrire un entier dans un système de numération donné, convertir entre bases, et effectuer des opérations dans une base donnée.
  \item Connaître et appliquer le petit théorème de Fermat pour accélérer un calcul de reste modulo un nombre premier.
  \item Traiter les exercices type examen national : résolution d'équations diophantiennes, systèmes de congruences.
\end{itemize}
\end{tcolorbox}

\vspace{4mm}

% ------------------- SECTION 1 -------------------
\section{Divisibilité dans $\mathbb{Z}$}

\begin{definition}[Divisibilité]
Soient $a,b\in\mathbb{Z}$. On dit que $b$ \textbf{divise} $a$ (ou que $a$ est un \textbf{multiple} de $b$, ou que $b$ est un \textbf{diviseur} de $a$), noté $b\mid a$, lorsqu'il existe $k\in\mathbb{Z}$ tel que $a=kb$.
\end{definition}

\begin{propriete}[Propriétés de la divisibilité]
Pour $a,b,c\in\mathbb{Z}$ :
\begin{itemize}[leftmargin=6mm, itemsep=0.5mm]
  \item \textbf{Réflexivité} : $a\mid a$. \quad Tout entier divise $0$ ; $1$ divise tout entier.
  \item \textbf{Transitivité} : si $a\mid b$ et $b\mid c$, alors $a\mid c$.
  \item Si $a\mid b$ et $a\mid c$, alors $a\mid(ub+vc)$ pour tous $u,v\in\mathbb{Z}$ (\textbf{combinaison linéaire}).
  \item Si $a\mid b$ et $b\mid a$ (avec $a,b\ne0$), alors $a=\pm b$.
  \item Si $a\mid b$ et $a,b\ne0$, alors $|a|\leq|b|$.
\end{itemize}
\end{propriete}

\begin{demo}
\emph{Transitivité} : $a\mid b\Rightarrow b=ka$ ; $b\mid c\Rightarrow c=k'b=k'ka$ : donc $a\mid c$ avec le coefficient $kk'$.

\smallskip
\emph{Combinaison linéaire} : $a\mid b\Rightarrow b=ka$ ; $a\mid c\Rightarrow c=k'a$. Alors $ub+vc=uka+vk'a=(uk+vk')a$ : $a$ divise bien $ub+vc$.

\smallskip
\emph{$a\mid b$ et $b\mid a\Rightarrow a=\pm b$} : $b=ka$ et $a=k'b=k'ka$, donc (si $a\ne0$) $kk'=1$, soit $k=k'=1$ ou $k=k'=-1$ (entiers de produit $1$) : $b=a$ ou $b=-a$.
\end{demo}

% ------------------- SECTION 2 -------------------
\section{Division euclidienne}

\begin{propriete}[Théorème de la division euclidienne]
Soient $a\in\mathbb{Z}$, $b\in\mathbb{N}^*$. Il existe un \textbf{unique} couple $(q,r)\in\mathbb{Z}\times\mathbb{N}$ tel que
\[
  a = bq+r, \qquad 0\leq r<b.
\]
$q$ est le \textbf{quotient}, $r$ le \textbf{reste} de la division euclidienne de $a$ par $b$.
\end{propriete}

\begin{demo}[Idée de la preuve]
\emph{Existence} : considérons l'ensemble $E=\{a-bk : k\in\mathbb{Z},\ a-bk\geq0\}$, non vide (pour $k$ assez négatif ou positif selon le signe de $a$). $E\subset\mathbb{N}$ possède donc un plus petit élément $r=a-bq$. Montrons $r<b$ : sinon $r-b=a-b(q+1)\geq0$ appartiendrait à $E$ et serait strictement plus petit que $r$, contredisant la minimalité de $r$.

\smallskip
\emph{Unicité} : si $a=bq+r=bq'+r'$ avec $0\leq r,r'<b$, alors $b(q-q')=r'-r$, avec $|r'-r|<b$. Comme $b\mid(r'-r)$ et $|r'-r|<b$, nécessairement $r'-r=0$, donc $r=r'$ puis $q=q'$.
\end{demo}

\begin{exemple}[Division euclidienne]
$29=6\times4+5$ (quotient $4$, reste $5$) ; $-17=6\times(-3)+1$ (attention : quotient $-3$, reste $1$, \textbf{pas} $-17=6\times(-2)-5$ car le reste doit être positif).
\end{exemple}

% ------------------- SECTION 3 -------------------
\section{Systèmes de numération}

\begin{definition}[Écriture d'un entier en base $b$]
Soit $b\in\mathbb{N}$, $b\geq2$. Tout entier $n\in\mathbb{N}^*$ s'écrit de façon \textbf{unique} sous la forme
\[
  n = a_p\,b^p + a_{p-1}\,b^{p-1} + \cdots + a_1\,b + a_0, \qquad
  a_p\ne0,\ \ 0\leq a_i\leq b-1 \text{ pour tout } i.
\]
On note alors $n = \overline{a_p a_{p-1}\cdots a_1 a_0}^{\,(b)}$, l'\textbf{écriture de $n$ en base $b$} (les $a_i$ sont les \textbf{chiffres} de $n$ en base $b$). Pour $b=10$ (base décimale usuelle), on omet en général l'exposant.
\end{definition}

\begin{propriete}[Obtention des chiffres par divisions euclidiennes successives]
Les chiffres $a_0,a_1,\ldots,a_p$ de l'écriture de $n$ en base $b$ s'obtiennent en effectuant des \textbf{divisions euclidiennes successives par $b$} : $a_0$ est le reste de $n$ par $b$, $a_1$ est le reste du quotient obtenu par $b$, et ainsi de suite jusqu'à un quotient nul.
\end{propriete}

\begin{exemple}[Conversion décimale $\to$ base $b$]
Écrivons $n=45$ en base $2$ (binaire).
\[
  45 = 2\times22+1, \quad 22=2\times11+0, \quad 11=2\times5+1, \quad 5=2\times2+1, \quad 2=2\times1+0, \quad 1=2\times0+1.
\]
En lisant les restes du dernier au premier : $45=\overline{101101}^{(2)}$.

\smallskip
\emph{Vérification} : $32+0+8+4+0+1=45$. $\checkmark$ (car $\overline{101101}^{(2)}=1\times2^5+0\times2^4+1\times2^3+1\times2^2+0\times2+1$).
\end{exemple}

\begin{propriete}[Conversion base $b$ $\to$ décimale]
Pour convertir $\overline{a_pa_{p-1}\cdots a_0}^{(b)}$ en base $10$, on calcule directement $\displaystyle\sum_{i=0}^pa_i\,b^i$.
\end{propriete}

\begin{exemple}[Conversion base $b$ $\to$ décimale, et opérations en base $b$]
Le nombre $\overline{2301}^{(4)}$ vaut, en base $10$ :
\[
  2\times4^3+3\times4^2+0\times4+1 = 128+48+0+1 = 177.
\]

\smallskip
\emph{Addition en base $b$} : pour additionner directement en base $b$ (sans repasser par la base $10$), on additionne chiffre par chiffre en \textbf{retenant} dès qu'une somme atteint $b$ (exactement comme la retenue en base $10$ dès qu'on atteint $10$). Par exemple, en base $5$ : $\overline{43}^{(5)}+\overline{24}^{(5)}$ : unités $3+4=7=1\times5+2$, on pose $2$ et retient $1$ ; dizaines $4+2+1=7=1\times5+2$, on pose $2$ et retient $1$. Résultat : $\overline{122}^{(5)}$ (vérification en base $10$ : $23+14=37=\overline{122}^{(5)}$ car $1\times25+2\times5+2=37$ $\checkmark$).
\end{exemple}

\begin{remarque}[Comparer deux nombres en base $b$]
Pour comparer deux entiers écrits dans la \textbf{même base} $b$, on compare d'abord leur nombre de chiffres (plus de chiffres $\Rightarrow$ plus grand), puis, à nombre de chiffres égal, on compare chiffre par chiffre depuis la gauche (poids fort), exactement comme en base $10$.
\end{remarque}

% ------------------- SECTION 4 -------------------
\section{PGCD, PPCM, algorithme d'Euclide}

\begin{definition}[PGCD]
Pour $a,b\in\mathbb{Z}$ non tous deux nuls, le \textbf{PGCD} (plus grand commun diviseur) de $a$ et $b$, noté $a\wedge b$ ou $\text{pgcd}(a,b)$, est le plus grand entier naturel qui divise à la fois $a$ et $b$.
\end{definition}

\begin{propriete}[Propriété fondamentale de l'algorithme d'Euclide]
Pour $a\in\mathbb{Z}$, $b\in\mathbb{N}^*$, si $a=bq+r$ (division euclidienne), alors
\[
  a\wedge b = b\wedge r.
\]
\end{propriete}

\begin{demo}
Montrons que l'ensemble des diviseurs communs de $a$ et $b$ coïncide avec celui de $b$ et $r$. Si $d\mid a$ et $d\mid b$, alors $d\mid(a-bq)=r$ : $d$ est un diviseur commun de $b$ et $r$. Réciproquement, si $d\mid b$ et $d\mid r$, alors $d\mid(bq+r)=a$ : $d$ est un diviseur commun de $a$ et $b$. Les deux paires ont donc exactement les mêmes diviseurs communs, donc le même plus grand.
\end{demo}

\begin{propriete}[Algorithme d'Euclide]
Pour calculer $a\wedge b$ : on effectue des divisions euclidiennes successives ($a$ par $b$, puis $b$ par le reste, etc.) jusqu'à obtenir un reste nul. Le \textbf{dernier reste non nul} est le PGCD.
\end{propriete}

\begin{exemple}[Calcul par l'algorithme d'Euclide]
Calculons $252\wedge108$.
\[
  252 = 108\times2+36, \qquad 108 = 36\times3+0.
\]
Le dernier reste non nul est $36$ : $\boxed{252\wedge108=36}$.
\end{exemple}

\begin{definition}[PPCM]
Le \textbf{PPCM} (plus petit commun multiple) de $a,b\in\mathbb{Z}^*$, noté $a\vee b$, est le plus petit entier naturel non nul multiple à la fois de $a$ et de $b$.
\end{definition}

\begin{propriete}[Relation PGCD-PPCM]
Pour $a,b\in\mathbb{N}^*$ :
\[
  (a\wedge b)\times(a\vee b) = a\times b.
\]
\end{propriete}

\begin{exemple}[Calcul de PPCM]
Avec $a=252$, $b=108$, $a\wedge b=36$ (calculé ci-dessus) :
\[
  a\vee b = \frac{252\times108}{36} = 252\times3 = 756.
\]
\end{exemple}

% ------------------- SECTION 5 -------------------
\section{Théorème de Bézout et théorème de Gauss}

\begin{definition}[Entiers premiers entre eux]
$a,b\in\mathbb{Z}$ sont \textbf{premiers entre eux} lorsque $a\wedge b=1$.
\end{definition}

\begin{propriete}[Théorème de Bézout]
Pour $a,b\in\mathbb{Z}$ non tous deux nuls, il existe $u,v\in\mathbb{Z}$ tels que
\[
  au+bv = a\wedge b.
\]
\textbf{Cas particulier (Bézout)} : $a$ et $b$ sont premiers entre eux si et seulement s'il existe $u,v\in\mathbb{Z}$ tels que $au+bv=1$.
\end{propriete}

\begin{remarque}[Algorithme d'Euclide étendu]
Les coefficients $u,v$ (identité de Bézout) s'obtiennent en \textbf{remontant} les étapes de l'algorithme d'Euclide.
\end{remarque}

\begin{exemple}[Détermination des coefficients de Bézout]
Trouvons $u,v\in\mathbb{Z}$ tels que $17u+5v=1$.

\smallskip
\textbf{Euclide} : $17=5\times3+2$ ; $5=2\times2+1$ ; $2=1\times2+0$.

\smallskip
\textbf{Remontée} : $1=5-2\times2$. Or $2=17-5\times3$, donc
\[
  1 = 5-(17-5\times3)\times2 = 5-17\times2+5\times6 = 5\times7-17\times2.
\]
\textbf{Conclusion} : $u=-2$, $v=7$ (vérification : $17\times(-2)+5\times7=-34+35=1$ $\checkmark$).
\end{exemple}

\begin{propriete}[Théorème de Gauss]
Si $a\mid bc$ et $a\wedge b=1$, alors $a\mid c$.
\end{propriete}

\begin{demo}
Par Bézout, comme $a\wedge b=1$, il existe $u,v$ tels que $au+bv=1$. En multipliant par $c$ : $auc+bvc=c$. Or $a\mid auc$ (facteur $a$) et $a\mid bvc$ (car $a\mid bc$ par hypothèse, donc $a\mid bcv$). Donc $a$ divise la somme $auc+bvc=c$.
\end{demo}

\begin{propriete}[Conséquences du théorème de Gauss]
\begin{itemize}[leftmargin=6mm, itemsep=0.5mm]
  \item Si $a\mid c$, $b\mid c$ et $a\wedge b=1$, alors $ab\mid c$.
  \item Si $p$ est premier et $p\mid ab$, alors $p\mid a$ ou $p\mid b$ (lemme d'Euclide, cas particulier de Gauss).
\end{itemize}
\end{propriete}

\begin{exemple}[Résoudre une équation diophantienne $ax+by=c$]
Résolvons dans $\mathbb{Z}^2$ : $17x+5y=3$.

\smallskip
\textbf{Solution particulière} : d'après l'exemple précédent, $17\times(-2)+5\times7=1$, donc en multipliant par $3$ : $17\times(-6)+5\times21=3$. Une solution particulière est $(x_0,y_0)=(-6,21)$.

\smallskip
\textbf{Solution générale} : si $(x,y)$ est une autre solution, alors $17(x-x_0)=-5(y-y_0)$, soit $17(x-x_0)=5(y_0-y)$. Comme $17\wedge5=1$, par Gauss, $17\mid(y_0-y)$, donc $y_0-y=17k$, $k\in\mathbb{Z}$, d'où $y=21-17k$ et $x=x_0+5k=-6+5k$.

\smallskip
\textbf{Conclusion} : $S=\{(-6+5k\,;21-17k) : k\in\mathbb{Z}\}$.
\end{exemple}

% ------------------- SECTION 6 -------------------
\section{Nombres premiers}

\begin{definition}[Nombre premier]
$p\in\mathbb{N}$, $p\geq2$, est \textbf{premier} lorsque ses seuls diviseurs positifs sont $1$ et $p$.
\end{definition}

\begin{propriete}[Décomposition en facteurs premiers -- admis]
Tout entier $n\geq2$ se décompose de façon \textbf{unique} (à l'ordre près) en produit de nombres premiers :
\[
  n = p_1^{\alpha_1}p_2^{\alpha_2}\cdots p_k^{\alpha_k}.
\]
\end{propriete}

\begin{propriete}[Infinité des nombres premiers -- Euclide]
L'ensemble des nombres premiers est \textbf{infini}.
\end{propriete}

\begin{demo}
Par l'absurde : supposons qu'il n'existe qu'un nombre fini de premiers $p_1,\ldots,p_k$. Posons $N=p_1p_2\cdots p_k+1$. $N$ n'est divisible par aucun des $p_i$ (le reste de la division de $N$ par $p_i$ vaut $1$). Or $N\geq2$ admet un diviseur premier $p$ (par la décomposition en facteurs premiers) : $p$ ne peut être aucun des $p_i$, contredisant l'hypothèse qu'il n'y en avait que $k$.
\end{demo}

% ------------------- SECTION 7 -------------------
\section{Congruences modulo $n$}

\begin{definition}[Congruence]
Pour $n\in\mathbb{N}^*$, $a,b\in\mathbb{Z}$, on dit que $a$ est \textbf{congru à $b$ modulo $n$}, noté $a\equiv b\ [n]$, lorsque $n\mid(a-b)$, c'est-à-dire lorsque $a$ et $b$ ont le \textbf{même reste} dans la division euclidienne par $n$.
\end{definition}

\begin{propriete}[Propriétés des congruences]
La relation $\equiv[n]$ est \textbf{réflexive, symétrique, transitive}. De plus, si $a\equiv b\,[n]$ et $c\equiv d\,[n]$ :
\[
  a+c \equiv b+d\ [n], \qquad ac \equiv bd\ [n], \qquad a^k \equiv b^k\ [n]\ (k\in\mathbb{N}).
\]
\end{propriete}

\begin{demo}
\emph{Compatibilité avec le produit} : $a\equiv b\,[n]\Rightarrow a-b=kn$ ; $c\equiv d\,[n]\Rightarrow c-d=k'n$. Alors $ac-bd=ac-bc+bc-bd=c(a-b)+b(c-d)=ckn+bk'n=n(ck+bk')$, donc $n\mid(ac-bd)$, soit $ac\equiv bd\,[n]$.
\end{demo}

\begin{exemple}[Critère de divisibilité par $9$]
Montrons qu'un entier est divisible par $9$ si et seulement si la somme de ses chiffres l'est.

\smallskip
Comme $10\equiv1\,[9]$, on a $10^k\equiv1^k=1\,[9]$ pour tout $k$. Si $N=\overline{a_pa_{p-1}\ldots a_1a_0}=\displaystyle\sum_{i=0}^pa_i10^i$, alors
\[
  N \equiv \sum_{i=0}^pa_i\times1 = \sum_{i=0}^pa_i \ [9].
\]
Donc $N\equiv0\,[9] \iff \sum a_i\equiv0\,[9]$, ce qui est exactement le critère annoncé.
\end{exemple}

\begin{exemple}[Résoudre une équation de congruence]
Résolvons $3x\equiv2\,[7]$ dans $\mathbb{Z}$.

\smallskip
Comme $3\wedge7=1$, $3$ est inversible modulo $7$. On cherche $u$ tel que $3u\equiv1\,[7]$ : $3\times5=15\equiv1\,[7]$ ($15=2\times7+1$). Donc l'inverse de $3$ modulo $7$ est $5$.

\smallskip
En multipliant l'équation par $5$ : $15x\equiv10\,[7]$, soit $x\equiv10\equiv3\,[7]$ (car $15\equiv1\,[7]$).

\smallskip
\textbf{Conclusion} : $S=\{3+7k : k\in\mathbb{Z}\}$.
\end{exemple}

\begin{exemple}[Calcul de reste par puissance modulaire]
Déterminons le reste de $2^{100}$ dans la division euclidienne par $7$.

\smallskip
Puissances de $2$ modulo $7$ : $2^1\equiv2$, $2^2\equiv4$, $2^3\equiv1\,[7]$ (car $8=7+1$). L'ordre de $2$ modulo $7$ est donc $3$ (période).

\smallskip
$100=3\times33+1$, donc $2^{100}=(2^3)^{33}\times2^1\equiv1^{33}\times2=2\,[7]$.

\smallskip
\textbf{Conclusion} : le reste de $2^{100}$ par $7$ est $2$.
\end{exemple}

\begin{propriete}[Petit théorème de Fermat]
Soit $p$ un nombre \textbf{premier} et $a\in\mathbb{Z}$ tel que $p\nmid a$ (c'est-à-dire $a\wedge p=1$). Alors
\[
  a^{p-1} \equiv 1 \ [p].
\]
Sous une forme équivalente, valable pour \textbf{tout} $a\in\mathbb{Z}$ (y compris si $p\mid a$) :
\[
  a^p \equiv a \ [p].
\]
\end{propriete}

\begin{demo}[Idée de la preuve -- admis dans le détail]
On admet ce résultat au niveau du bac (sa démonstration standard utilise le fait que les restes non nuls modulo $p$, multipliés par $a$, redonnent une permutation des mêmes restes, puis on simplifie un produit). L'important est de savoir \textbf{l'utiliser} pour accélérer un calcul de reste modulo un nombre premier.
\end{demo}

\begin{exemple}[Application directe de Fermat]
Déterminons le reste de $2^{100}$ modulo $7$ en utilisant le petit théorème de Fermat (à comparer avec la méthode de l'exemple précédent).

\smallskip
$7$ est premier et $7\nmid2$, donc par Fermat : $2^{6}\equiv1\,[7]$.

\smallskip
$100 = 6\times16+4$, donc $2^{100}=\left(2^6\right)^{16}\times2^4\equiv1^{16}\times2^4=16\equiv2\,[7]$ (car $16=2\times7+2$).

\smallskip
\textbf{Conclusion} : on retrouve bien un reste de $2$, obtenu ici directement à partir de l'ordre $p-1=6$ donné par Fermat, sans avoir à chercher la période « à la main ».
\end{exemple}

\begin{remarque}[Quand utiliser Fermat plutôt que chercher la période soi-même]
Le petit théorème de Fermat donne \textbf{immédiatement} un exposant $p-1$ qui fonctionne (dès que $p$ est premier et $p\nmid a$), sans avoir à calculer les premières puissances pour repérer la période. Attention : la \textbf{période réelle} peut être un diviseur strict de $p-1$ (ex. ci-dessus, la période de $2$ modulo $7$ est $3$, diviseur de $6=p-1$) ; les deux méthodes sont donc valables, Fermat étant plus rapide à mettre en pratique sans calcul préalable.
\end{remarque}

% ------------------- SECTION 8 -------------------
\section{Exercices corrigés -- niveau examen national SM}

\begin{exemple}[Exercice 1 -- diviseurs communs et PGCD]
Déterminer tous les couples $(x,y)\in(\mathbb{N}^*)^2$ tels que $x\wedge y=6$ et $x+y=36$.

\smallskip
\textbf{Corrigé.} Posons $x=6a$, $y=6b$ avec $a\wedge b=1$ (propriété de la division par le PGCD). Alors $6a+6b=36\iff a+b=6$.

\smallskip
Couples $(a,b)$ avec $a+b=6$ et $a\wedge b=1$ : $(1,5)$, $(5,1)$ conviennent ($1\wedge5=1$) ; $(2,4)$ ne convient pas ($2\wedge4=2$) ; $(3,3)$ ne convient pas ($3\wedge3=3$).

\smallskip
\textbf{Conclusion} : $(x,y)\in\{(6,30)\,;(30,6)\}$.
\end{exemple}

\begin{exemple}[Exercice 2 -- équation diophantienne avec contrainte]
Une boutique vend des stylos à $12$ DH et des cahiers à $17$ DH. Un client dépense exactement $200$ DH. Combien de stylos et de cahiers a-t-il pu acheter ?

\smallskip
\textbf{Corrigé.} On cherche $(x,y)\in\mathbb{N}^2$ tels que $12x+17y=200$.

\smallskip
$12\wedge17=1$ (consécutifs dans l'algorithme, ou vérification directe), donc l'équation a des solutions entières. Par tâtonnement/Bézout : $12\times(-7)+17\times5=-84+85=1$, donc $12\times(-1400)+17\times1000=200$ est une solution particulière (peu pratique) ; cherchons directement par essais bornés : $y$ doit vérifier $17y\leq200$, soit $y\leq11$, et $200-17y\equiv0\,[12]$.

\smallskip
En testant : $y=4\Rightarrow200-68=132=12\times11$. Solution $(x,y)=(11,4)$.

\smallskip
\textbf{Conclusion} : le client a acheté $11$ stylos et $4$ cahiers (solution dans $\mathbb{N}^2$ ; la solution générale entière est $(11-17k\,;4+12k)$, dont seul $k=0$ donne des valeurs positives raisonnables ici).
\end{exemple}

\begin{exemple}[Exercice 3 -- démonstration avec Gauss]
Montrer que pour tout $n\in\mathbb{Z}$, $n(n+1)(2n+1)$ est divisible par $6$.

\smallskip
\textbf{Corrigé.} \emph{Divisibilité par $2$} : parmi $n$ et $n+1$, l'un est pair : $2\mid n(n+1)$, donc $2\mid n(n+1)(2n+1)$.

\smallskip
\emph{Divisibilité par $3$} : étudions les restes modulo $3$. Si $n\equiv0\,[3]$ : $3\mid n$. Si $n\equiv1\,[3]$ : $2n+1\equiv3\equiv0\,[3]$. Si $n\equiv2\,[3]$ : $n+1\equiv3\equiv0\,[3]$. Dans tous les cas, $3$ divise l'un des trois facteurs.

\smallskip
Comme $2\wedge3=1$ et $2\mid N$, $3\mid N$ (où $N=n(n+1)(2n+1)$), par la conséquence du théorème de Gauss : $6=2\times3\mid N$.
\end{exemple}

\begin{exemple}[Exercice 4 -- système de congruences]
Résoudre le système $\begin{cases} x\equiv2\,[5] \\ x\equiv3\,[7] \end{cases}$.

\smallskip
\textbf{Corrigé.} $x\equiv2\,[5]\Rightarrow x=5k+2$. En substituant dans la seconde : $5k+2\equiv3\,[7]\iff5k\equiv1\,[7]$.

\smallskip
Inverse de $5$ modulo $7$ : $5\times3=15\equiv1\,[7]$, donc l'inverse est $3$. D'où $k\equiv3\,[7]$, soit $k=7m+3$.

\smallskip
Alors $x=5(7m+3)+2=35m+17$.

\smallskip
\textbf{Conclusion} : $S=\{17+35m : m\in\mathbb{Z}\}$, soit $x\equiv17\,[35]$ (théorème des restes chinois, $5\times7=35$).
\end{exemple}

\begin{exemple}[Exercice 5 -- reste par Fermat]
Déterminer le reste de la division euclidienne de $7^{2024}$ par $13$.

\smallskip
\textbf{Corrigé.} $13$ est premier et $13\nmid7$ : par Fermat, $7^{12}\equiv1\,[13]$.

\smallskip
$2024 = 12\times168+8$, donc $7^{2024}=\left(7^{12}\right)^{168}\times7^8\equiv7^8\,[13]$.

\smallskip
Calculons $7^8$ modulo $13$ par carrés successifs : $7^2=49\equiv10\,[13]$ (car $49=3\times13+10$) ; $7^4\equiv10^2=100\equiv9\,[13]$ (car $100=7\times13+9$) ; $7^8\equiv9^2=81\equiv3\,[13]$ (car $81=6\times13+3$).

\smallskip
\textbf{Conclusion} : le reste de $7^{2024}$ par $13$ est $3$.
\end{exemple}

\begin{exemple}[Exercice 6 -- système de numération et divisibilité]
Soit $n=223$ (en base $10$). Écrire $n$ en base $7$, puis dire si $n$ est divisible par $3$.

\smallskip
\textbf{Corrigé.} \emph{Conversion en base $7$} : divisions euclidiennes successives par $7$.
\[
  223=7\times31+6, \qquad 31=7\times4+3, \qquad 4=7\times0+4.
\]
En lisant les restes du dernier quotient au premier reste obtenu : $n=\overline{436}^{(7)}$.

\smallskip
\emph{Vérification} : $4\times7^2+3\times7+6 = 196+21+6 = 223$. $\checkmark$

\smallskip
\emph{Divisibilité par $3$} : $223=3\times74+1$, donc $223\equiv1\,[3]$ : $n$ n'est \textbf{pas} divisible par $3$.
\end{exemple}

% ------------------- SECTION 9 -------------------
\section{Méthodes et pièges : la boîte à outils de l'examen}

\subsection{Ce qu'on te demande, ce que tu fais}

\begin{center}
\renewcommand{\arraystretch}{1.45}
\sffamily\small
\begin{tabular}{>{\raggedright\arraybackslash}p{7.2cm} >{\raggedright\arraybackslash}p{8cm}}
\rowcolor{qtealdark} \color{white}\bfseries Ce qu'on te demande & \color{white}\bfseries Ton réflexe \\
\rowcolor{tealpale} Calculer un PGCD & Algorithme d'Euclide (divisions successives), dernier reste non nul. \\
Trouver $u,v$ tels que $au+bv=d$ & Euclide pour obtenir $d$, puis remonter les calculs (Bézout). \\
\rowcolor{tealpale} Résoudre $ax+by=c$ dans $\mathbb{Z}^2$ & Vérifier $a\wedge b\mid c$, trouver une solution particulière (Bézout), écrire la solution générale via Gauss. \\
Prouver $a\wedge b=1$ & Bézout (trouver $u,v$) ou Euclide (dernier reste $=1$). \\
\rowcolor{tealpale} Démontrer une divisibilité par un produit ($6$, $12$...) & Décomposer en facteurs premiers entre eux, prouver chaque divisibilité séparément, conclure par Gauss. \\
Calculer un reste de grande puissance & Chercher la période des puissances modulo $n$, réduire l'exposant. \\
\rowcolor{tealpale} Résoudre $ax\equiv b\,[n]$ & Trouver l'inverse de $a$ modulo $n$ (Bézout), multiplier les deux membres. \\
Critère de divisibilité (par $9$, $11$...) & Utiliser $10\equiv1\,[9]$ ou $10\equiv-1\,[11]$ et sommer/alterner les chiffres. \\
\rowcolor{tealpale} Calculer un reste de grande puissance modulo un nombre \textbf{premier} $p$ & Petit théorème de Fermat : $a^{p-1}\equiv1\,[p]$ si $p\nmid a$, réduire l'exposant modulo $p-1$. \\
Convertir un entier en base $b$ (ou l'inverse) & Divisions euclidiennes successives par $b$ (lire les restes du dernier au premier) ; ou $\sum a_ib^i$ pour l'inverse. \\
\end{tabular}
\end{center}

\subsection{Les pièges qui coûtent des points}

\begin{remarque}[Erreurs relevées chaque année par les correcteurs]
\begin{itemize}[leftmargin=6mm, itemsep=0.5mm]
  \item Dans la division euclidienne d'un nombre \textbf{négatif}, oublier que le reste doit rester dans $[0\,;b[$ (pas de reste négatif).
  \item Confondre $a\wedge b$ (PGCD) et $a\vee b$ (PPCM) dans une formule.
  \item Appliquer le théorème de Gauss sans vérifier l'hypothèse $a\wedge b=1$.
  \item Oublier de vérifier que $a\wedge b$ divise $c$ avant de chercher à résoudre $ax+by=c$ (condition nécessaire d'existence de solutions).
  \item Écrire une solution générale d'équation diophantienne sans avoir d'abord soustrait la solution particulière (étape de la preuve par Gauss).
  \item Oublier de vérifier les coefficients de Bézout trouvés (toujours vérifier $au+bv=d$ numériquement).
  \item Appliquer le petit théorème de Fermat avec un $p$ \textbf{non premier}, ou sans vérifier $p\nmid a$.
  \item Dans une conversion de base, oublier qu'un chiffre $a_i$ doit vérifier $0\leq a_i\leq b-1$ (jamais un chiffre $\geq b$).
\end{itemize}
\end{remarque}

% ------------------- RESUME -------------------
\vspace{4mm}
\begin{tcolorbox}[
  enhanced, boxrule=0pt, arc=2mm,
  interior style={left color=qtealdark, right color=qteal},
  left=6mm, right=6mm, top=4mm, bottom=4mm,
  fontupper=\color{white}\sffamily,
  title={\sffamily\bfseries\color{white}\ding{72}~~À retenir},
  colbacktitle=qtealdark, coltitle=white, boxrule=0pt
]
\begin{itemize}[leftmargin=6mm, label={\color{qamber}\ding{212}}, itemsep=1mm]
  \item Division euclidienne : $a=bq+r$, $0\leq r<b$, unique.
  \item Euclide : $a\wedge b=b\wedge r$ ; calcul par divisions successives, dernier reste non nul.
  \item Bézout : $\exists u,v,\ au+bv=a\wedge b$ ; $a\wedge b=1\iff\exists u,v,\ au+bv=1$.
  \item Gauss : $a\mid bc$ et $a\wedge b=1\Rightarrow a\mid c$ ; conséquence : $a\mid c$, $b\mid c$, $a\wedge b=1\Rightarrow ab\mid c$.
  \item Équation $ax+by=c$ : solutions ssi $(a\wedge b)\mid c$ ; solution générale via Gauss à partir d'une solution particulière.
  \item Congruences : $a\equiv b\,[n]$ compatible avec $+,\times,$ puissances ; utile pour restes et critères de divisibilité.
  \item Fermat : $p$ premier, $p\nmid a$ $\Rightarrow a^{p-1}\equiv1\,[p]$ (et $a^p\equiv a\,[p]$ pour tout $a$).
  \item Base $b$ : $n=\sum a_ib^i$ avec $0\leq a_i<b$ ; chiffres obtenus par divisions euclidiennes successives par $b$.
\end{itemize}
\end{tcolorbox}

\end{document}
